% Compile with XeLaTeX or LuaLaTeX on an updated TeX Live 2024 or later.
\documentclass[11pt,a4paper]{article}
\usepackage[margin=25mm]{geometry}
\usepackage[language=assamese]{latexbangla}
\setbanglamainfont{Kalpurush}
\setmainfont{TeX Gyre Termes} % Optional: change the English font.
\usepackage{amsmath,amsthm}
\usepackage[hidelinks]{hyperref}
\listfiles

% Document-specific environments belong in the document or its class.
\newtheorem{theorem}{উপপাদ্য}[section]

\begin{document}
\section{পাইথাগোৰাছৰ উপপাদ্য}
এখন আয়তাকাৰ বাগিচাৰ দৈৰ্ঘ্য ৪ মিটাৰ আৰু প্ৰস্থ ৩ মিটাৰ। বাগিচাখনৰ এটা কোণৰ
পৰা বিপৰীত কোণলৈ পোনে পোনে যাবলৈ কিমান দূৰ খোজ কাঢ়িব লাগিব?\footnote{ইয়াত বাগিচাখন
সমতল বুলি ধৰা হৈছে আৰু সকলো দৈৰ্ঘ্য একে এককত জোখা হৈছে।}
উত্তৰ পাবলৈ আমি \textit{Pythagorean theorem} ব্যৱহাৰ কৰিব পাৰোঁ।

The shortest route across the garden is its diagonal—অসমীয়াত যাক কৰ্ণ বুলি কওঁ।
Imagine drawing that route on a plan: the two sides of the garden and the
diagonal form a right triangle. With sides of 3 and 4 metres, walking around
the edge takes 7 metres. But can we cross the garden in just 5 metres?
\textbf{A sketch suggests the answer; নিশ্চিত হ'বলৈ এটা প্ৰমাণ লাগে।}

\begin{theorem}
সমকোণী ত্ৰিভুজৰ লম্ব বাহু দুটাৰ দৈৰ্ঘ্য $a$ আৰু $b$ আৰু অতিভুজৰ দৈৰ্ঘ্য $c$ হ'লে,
\begin{equation}\label{eq:pythagoras}
  a^2+b^2=c^2.
\end{equation}
\end{theorem}
এই উপপাদ্যৰ এটা জ্যামিতিক প্ৰমাণ ইউক্লিডৰ \textit{Elements}-ত পোৱা যায়~\cite{euclid}।

\subsection{বাগিচাৰ কৰ্ণ আৰু ক্ষেত্ৰফল}
\eqref{eq:pythagoras}ৰ পৰা বাগিচাখনৰ কৰ্ণৰ দৈৰ্ঘ্য পাওঁ,
\begin{align*}
  c &= \sqrt{৩^{২}+৪^{২}}=\sqrt{২৫}=৫\text{ মিটাৰ},\\
  \text{বাগিচাৰ ক্ষেত্ৰফল} &= \text{দৈৰ্ঘ্য}\times\text{প্ৰস্থ}
    =৪\times৩=১২\text{ বৰ্গমিটাৰ।}
\end{align*}
\textbf{মনত ৰাখোঁ: উপপাদ্যটো কেৱল সমকোণী ত্ৰিভুজৰ ক্ষেত্ৰতহে প্ৰযোজ্য।}
বাগিচাখনৰ কাষেৰে ঘূৰি খোজ কঢ়াৰ সলনি কৰ্ণৰ ওপৰেৰে গ'লে সেয়ে ২ মিটাৰ পথ কমি যায়।
\textit{A shorter route, the same destination—একে গন্তব্যলৈ কম পথ।}

\autoref{tab:triples}ত এনে আৰু কেইটামান বাহুৰ দৈৰ্ঘ্য দিয়া হ'ল।
\begin{table}[ht]
  \centering
  \caption{সমকোণী ত্ৰিভুজৰ পূৰ্ণসংখ্যক বাহুৰ দৈৰ্ঘ্য}\label{tab:triples}
  \begin{tabular}{ccc}
    \hline
    প্ৰথম লম্ব বাহু & দ্বিতীয় লম্ব বাহু & অতিভুজ \\
    \hline
    ৩ & ৪ & ৫ \\
    ৫ & ১২ & ১৩ \\
    ৮ & ১৫ & ১৭ \\
    \hline
  \end{tabular}
\end{table}

\subsection{নিজে কৰোঁ}
\begin{enumerate}
  \item লম্ব বাহু দুটাৰ দৈৰ্ঘ্য ৬ আৰু ৮ একক হ'লে অতিভুজৰ দৈৰ্ঘ্য নিৰ্ণয় কৰা।
  \item এখন আয়তক্ষেত্ৰৰ দৈৰ্ঘ্য ১২ আৰু প্ৰস্থ ৫ একক।
    \begin{enumerate}
      \item আয়তক্ষেত্ৰখনৰ কৰ্ণৰ দৈৰ্ঘ্য কিমান?
      \item কৰ্ণৰ ওপৰেৰে কাটিলে প্ৰতিটো ত্ৰিভুজৰ ক্ষেত্ৰফল কিমান?
    \end{enumerate}
\end{enumerate}

\begin{thebibliography}{9}
  \bibitem{euclid} Euclid. \textit{Elements}, Book I,
  \href{https://mathcs.clarku.edu/~djoyce/elements/bookI/propI47.html}{Proposition 47}.
  Online edition by David E. Joyce.
\end{thebibliography}
\end{document}
